\documentclass[12pt]{article}
\usepackage[select language]{babel} 
\usepackage{packages/mygeneral}
\usepackage{packages/mycommands}
\usepackage{packages/tfg_layout} 
%\usepackage{packages/mytheorems_env}
%\usepackage{packages/dangbox} 
%\usepackage{packages/epibox} 
%\usepackage[backend=biber, style=numeric, sorting=none]{biblatex} 
%\addbibresource{bibliography.bib} 
\institution{
  name=UAM
} 
\pagelayout{
 rightheader=Física Atómica y Molecular,
 logo=false % set false/true to hide/show the lowo in the left header
} 

\mathtoolsset{showonlyrefs=true} % for showing number only in labeled equations
 \allowdisplaybreaks[2] in case you want your equations to be broken across pages
\let\vec\mathbf % display vectors in bold instead of with an arrow
\title{Método mejorado de obtención de términos L-S y j-j para electrones degenerados}
\author{Joan A. Mercado Tandazo \and Juan Manuel Sánchez Arrúa} 
\begin{document}
\maketitle
\hrule{\textwidth}
\resetnumpagetitle
\tableofcontents
\resetnumpagetitle
\pagestyle{academic} 
\section{Introducción} 
El método que se presenta a continuación es una mejora del método tradicional que se presenta en el curso de Física Atómica y Molecular, el cual suele estar basado en tablas y contabilidad de posibilidades según estas se le ocurran al estudiante. El método que presentamos reduce el esfuerzo de contabilidad gracias a la combinatoria y simetrías de los términos, sistematizando el proceso y siendo fácilmente extensible a un número arbitrario de electrones.\par
El procedimiento consiste en centrarse en el número cuántico $M_{L}$ (o $M_{J}$ para el método j-j) e ir obteniendo los términos de acuerdo al número de combinaciones posibles de los electrones para dar ese número cuántico. Conviene empezar desde el máximo valor posible de $M_{L}$ (o $M_{J}$) e ir reduciendo su valor hasta llegar a $M_{L}=0$ (o $M_{J}=0$) (hasta ese valor será suficiente gracias a las simetrías). Esta reducción se puede realizar sistemáticamente gracias a la aplicación del operador $L_{-}$ (o $J_{-}$) como veremos en los ejemplos. Según las distintas combinaciones que surjan, y teniendo en cuenta el spin para el caso de L-S, obtendremos los términos correspondientes, que deberemos ir anotando a medida que los vayamos encontrando, puesto que la siguiente reducción los necesitará en la contabilidad. 
\section{Términos L-S} 
Supongamos que tenemos $N$ electrones en un mismo subnivel $l$, y queremos obtener los términos L-S correspondientes. Es decir, una configuración electrónica de la forma $l^{N}$ donde asumimos que $N \leq 1/2(2l + 1)\times (2s+1)$ pues por la simetría ``electrón-hueco'' el caso de $N > 1/2(2l + 1)\times (2s+1)$ se puede obtener a partir del caso de $N' = 2(2l + 1)\times (2s+1) - N$ electrones. Por teoría del momento angular, sabemos que el número cuántico total de momento orbital $L$ puede tomar valores entre $N \times L$ y $0$ o $l$ según $N$ sea par o impar. Esto nos da una idea de qué valores de $L$ podemos esperar en los términos L-S. Sin emgargo, una parte fundamental del método es estudiar el spin. Este tomará valores entre $S = N \times s$ y $0$ o $1/2$ según $N$ sea par o impar. Como las posibilidades dela tercera componente del spin individual son menores a las del momento orbital (individual) entonces le asignamos a este un rol secundario. De esta forma tenemos el siguiente panorama: para cada combinación $(M_{L},M_{S})$ tenemos un conjunto de posibilidades $(m_{1}^{\sigma_{1}}, m_{2}^{\sigma_{2}},\dots m_{N}^{\sigma_{N}})$ donde $m_{i}$ es la tercera componente del momento orbital del electrón $i$ y $\sigma_{i}$ es la tercera componente del spin del electrón $i$. Además, por teoría del momento angular deben cumplir: 
\begin{align}
  \sum_{i=1}^{N} &m_{i} = M_{L} & \sum_{i=1}^{N} &\sigma_{i} = M_{S} \label{eq: suma de 3as componentes} 
\end{align}

Cada uno de estos $m_{i}^{\sigma_{i}}$ son distintos para cumplir el principio de exclusión de Pauli. Existen $2\times(2l + 1))$ posibilidades de $m_{i}^{\sigma_{i}}$ para cada electrón, y por lo tanto, el número de formas posibles de tomar $N$ electrones distintos de entre estas $2\times(2l + 1))$ posibilidades es:
\begin{align*}
\# \text{combinaciones} = \binom{2 \times (2l+1)}{N}
\end{align*}
Lo cual se puede comprobar que coincide con el resultado del método tradicional, pues de hecho, es una explicación de ese resultado. Esto también nos quiere decir que la degeneración acumulada de los términos L-S que obtendremos ha de coincidir con este número de combinaciones, lo cual será una prueba de verificación del método. \par 
A continuación centrémonos en $M_{L}$ y busquemos el número más alto posible que puede tener. Por \eqref{eq: suma de 3as componentes} tendremos que buscar los máximos valores posibles de $m_{i}$, no obstante respetando el principio de exclusión. Para ello establecemos la regla de que, \textbf{en la suma que escribamos de $m$'s  (partición) no podemos repetir el mismo $m$ 2 veces}. Una vez tenemos la partición habremos obtenido el máximo valor de $M_{L}$ que a su vez corresponderá al máximo valor de $L$. Entonces, nos queda averiguar el spin. Para ello, debereos sumar las $\sigma$'s que, por supuesto, no hemos llevado su contabilidad, porque no es necesario. En el caso de $M_{L}$ más alto, es sencillo averiguar $M_{S}$ porque tendremos muchas repeticiones de $m$'s lo cual significa que las terceras componentes de spin se anulan entre sí y solo tendremos que fijarnos en la $m$ no repetida para averiguar $M_{S}$. Si esa $m$ no existe entonces estaremos en el caso con $N$ par y $M_{S} = 0$, lo cual inducirá a que $S=0$ para ese término de máxima $L$. Y si $m$ no repetida existe entonces estaremos en el caso con $N$ impar y $M_{S} = 1/2$, lo cual inducirá a que $S=1/2$ para ese término de máxima $L$. De ahí podremos extraer el término L-S correspondiente con su degeneración, que será importante para la contabilidad de los siguientes términos. \par
A continuación, lo siguiente es aplicar el operador $L_{-}$ para reducir el valor de $M_{L}$. Si bien recordamos, dicho operador viene dado por: 
\begin{align*}
L_{-} = \sum_{i=1}^{N} l_{-}^{(i)}
\end{align*}
Y si aplicamos $L_{-}$ a un estado con $M_{L}$ dado, entonces el resultado será un estado con $M_{L} - 1$. Pero además: 
\begin{align*}
L_{-} \ket{M_{L},M_{S}} = \sum_{i=0}^{N} l_{-}^{(i)} \ket{m_{1},\sigma_{1}; m_{2},\sigma_{2}; \dots m_{N},\sigma_{N}} = \sum_{i=0}^{N} a_{i} \ket{\dots m_{i} - 1, \sigma_{i} \dots}
\end{align*}
Donde $a_{i}$ por construcción será tal que cumpla el principio de exclusión de pauli, es decir, $a_{i} = 0$ si $m_{i} - 1$ ya está repetido. Esto se traduce en lenguaje de la partición que hemos hecho a que a cada sumando de la partición podemos reducirlo en una cifra para obtener una partición de $M_{L}-1$ siempre y cuando esa reducción no implique repetir más de 2 veces un $m$. Es posible, y de hecho favorable a la hora de reducir el esfuerzo, que a la hora de aplicar $L_{-}$ se obtenga una partición de $M_{L}-1$ que ya hemos anotado. En ese caso, no la volvemos a repetir, simplemente la ignoramos. Así podremos generar todas las particiones de $M_{L}-1$ a partir de las particiones de $M_{L}$, y sistemáticamente llegar a $M_{L} = 0$. \par 
Para cada uno de los $M_{L}$ tendremos que buscar los posibles $M_{S}$ correspondientes. En ese caso tendremos que generalizar lo que vimos para el caso de $M_{L}$ máximo. Si 
\subsection{Ejemplo: 2 \texorpdfstring{$e^{-}$}{e^-}} 

\subsection{Ejemplo: 3 \texorpdfstring{$e^{-}$}{e^-}} 

\subsection{Ejemplo: 4 \texorpdfstring{$e^{-}$}{e^-}}

\section{Términos j-j}

\section{Ejemplo: 2 \texorpdfstring{$e^{-}$}{e^-}}

\section{Ejemplo: 3 \texorpdfstring{$e^{-}$}{e^-}}

\section{Ejemplo: 4 \texorpdfstring{$e^{-}$}{e^-}}
\end{document} 

